\chapter{ABI Classification and Foreign-Call Lowering}
\label{ch:ffi-abi}

Passing eight integer arguments exposes a boundary that a two-argument demonstration never reaches.
Some values enter registers and others enter the stack area.
Adding a floating argument exposes another register class; returning a large aggregate can add an invisible result pointer.
An FFI implements the selected calling convention, not merely a cast to a convenient function-pointer type.

The examples use the ordinary System V AMD64 ABI for scalar integers, doubles, and two particular structures.
They exclude vector extensions, packed structures, bit-fields, and the full aggregate classification algorithm.
The native experiment delegates call construction to libffi and compares results with C calls compiled from the same declarations.

\section{Classify before assigning locations}
\label{sec:ch70-classify}
The first six integer-class argument slots use \code{RDI}, \code{RSI}, \code{RDX}, \code{RCX}, \code{R8}, and \code{R9}.
Floating arguments use a separate SSE register sequence beginning at \code{XMM0}.
Eligible aggregates are classified in eight-byte units, and assignment can require more than one register class~\cite{matz-2012-amd64}.
Classification depends on field layout and type categories, not just total size.

Our \code{Pair} contains a double followed by a signed 32-bit tag.
On the tested ABI it has size $16$, alignment $8$, and offsets $0$ and $8$.
Its two units classify as SSE and INTEGER.
A three-element aggregate of unsigned 64-bit integers instead uses memory in the ordinary profile considered here.
Table~\ref{tab:ch70-types} records the distinctions checked by the test fixture.

\begin{table}[htbp]
\caption{Equal source-level argument counts can demand different native resources.}
\label{tab:ch70-types}
\centering
\begin{tabular}{lll}
\toprule
Type & Tested layout & Argument treatment \\
\midrule
signed 32-bit integer & $4$ bytes & integer-class location \\
double & $8$ bytes & SSE-class location \\
Pair: double, int32 & $16$ bytes, alignment $8$ & SSE plus INTEGER, or stack \\
Triple: three uint64 & $24$ bytes, alignment $8$ & memory argument \\
\bottomrule
\end{tabular}
\end{table}

A memory-class argument and a memory-class return do not mean the same thing.
An ordinary memory argument is passed in the caller's stack argument area under this ABI.
A memory-class return uses caller-provided result storage whose address is passed through a hidden parameter; this consumes an integer argument register.
Other ABIs make different choices, so generic IR should retain the distinction until target lowering.

\section{Reserve aggregate resources atomically}
\label{sec:ch70-registers}
An aggregate that needs several registers must not partially consume them and spill only the remainder when one class runs out.
Insufficient registers cause the entire eligible aggregate to be passed on the stack and tentative assignments to be rolled back~\cite{matz-2012-amd64}.
Algorithm~\ref{alg:ch70-reserve} implements the restricted class-count model in the Python companion.
It computes register-versus-stack decisions, not stack byte offsets.

\begin{algorithm}[htbp]
\caption{Assign registers in the restricted scalar and two-unit aggregate profile.}
\label{alg:ch70-reserve}
\begin{pseudocode}
\Require classified arguments; six integer and eight SSE slots
\Ensure all-or-stack assignment for each supported aggregate
\State initialize used-slot counters, including any hidden result pointer
\ForAll{arguments in declaration order}
  \If{memory-class or either register class lacks capacity}
    \State assign the whole argument to the stack argument area
  \Else
    \State reserve all required integer and SSE slots together
    \State assign the argument's units to those slots
  \EndIf
\EndFor
\end{pseudocode}
\end{algorithm}

Restricted assignment costs $O(A)$ for $A$ already classified arguments.
General classification recursively inspects aggregate types and applies target-specific exceptional cases; this count-only model does not implement it.
Stack placement additionally accounts for size, alignment, and the total outgoing area.
Immediately before an ordinary call instruction the stack is $16$-byte aligned in this profile; the pushed return address shifts the callee-entry position by eight bytes.

Table~\ref{tab:ch70-pressure} works six integer arguments followed by Pair and a double.
The six integers fill the integer bank.
Pair cannot obtain its integer component, so it goes entirely to the stack without consuming \code{XMM0}.
The final double therefore still uses \code{XMM0}.
Figure~\ref{fig:ch70-banks} visualizes this decision.

\begin{table}[htbp]
\caption{An aggregate spill leaves the unused floating bank available to later arguments.}
\label{tab:ch70-pressure}
\centering
\begin{tabular}{lll}
\toprule
Argument & Demand & Assignment \\
\midrule
first six uint64 values & one integer slot each & six integer registers \\
Pair & one integer, one SSE & entire Pair on stack \\
final double & one SSE slot & XMM0 \\
\bottomrule
\end{tabular}
\end{table}

\begin{figure}[htbp]
\centering
\begin{tikzpicture}[node distance=12mm and 12mm]
\node[basic block] (pair) {Pair requires both banks};
\node[basic block,below left=of pair] (gp) {integer bank: full};
\node[basic block,below right=of pair] (fp) {SSE bank: free};
\node[basic block,below=of gp] (stack) {whole Pair on stack};
\draw[cfg edge] (pair)--(gp); \draw[cfg edge] (pair)--(fp); \draw[cfg edge] (gp)--(stack);
\end{tikzpicture}
\caption{A free register in one bank does not permit a partial aggregate assignment.}
\label{fig:ch70-banks}
\end{figure}

\section{Give call metadata a lifetime}
\label{sec:ch70-libffi}
libffi's \code{ffi\_cif} describes the ABI, return type, and argument types.
The runtime supplies addresses of actual argument objects, not values packed into pointer-sized integers.
A pointer argument therefore needs the address of a pointer-valued slot; passing the pointee directly adds an indirection error.
Type descriptors and their referenced arrays must remain alive while the prepared interface uses them~\cite{libffi-2026-manual}.

The companion's \code{FxPlan} owns its argument-type array and contains an interface pointing into that array.
It must remain at a stable address after preparation; copying the structure would leave internal pointers referring to old storage.
The scalar API supports signed 32-bit integers, unsigned 64-bit integers, doubles, and pointers, with at most eight arguments.
Aggregate and variadic tests use libffi directly rather than exposing those features through the small scalar-plan API.

\begin{algorithm}[htbp]
\caption{Invoke a known native target through an address-stable prepared plan.}
\label{alg:ch70-call}
\begin{pseudocode}
\Require prepared plan, trusted matching function address, typed values, output
\Ensure ABI-conformant call or rejection before invocation
\State validate argument count and value kinds against the plan
\State build addresses of correctly typed argument slots
\State select aligned return storage of sufficient size
\State invoke the target through the prepared libffi interface
\State normalize the native return into the wrapper's value kind
\State publish the wrapper result
\end{pseudocode}
\end{algorithm}

For $A$ scalar arguments, wrapper setup is $O(A)$, excluding native execution and ABI-specific libffi work.
For integral return types narrower than the machine register size, libffi can write a widened result; the companion uses \code{ffi\_sarg} storage for its signed 32-bit return channel~\cite{libffi-2026-manual}.
Allocating only four bytes for that channel on a wider machine risks overwriting adjacent memory.
Arguments instead use their exact declared scalar object types.

Figure~\ref{fig:ch70-indirection} separates argument-value addresses from the pointed-to output.
The metadata cannot establish that a symbol actually has the supplied signature.
That remains a trusted declaration and library-identity assumption.

\begin{figure}[htbp]
\centering
\begin{tikzpicture}[node distance=12mm]
\node[basic block] (slots) {slots: a, b, output pointer};
\node[basic block,below left=of slots] (call) {libffi reads slot addresses};
\node[basic block,below right=of slots] (storage) {writable output integer};
\draw[cfg edge] (slots)--(call);
\draw[cfg edge] (slots)--node[right,font=\small]{pointer value}(storage);
\end{tikzpicture}
\caption{The pointer-valued slot refers to separate storage; libffi receives the slot's address.}
\label{fig:ch70-indirection}
\end{figure}

\section{Promote variadic arguments}
\label{sec:ch70-varargs}
C's default promotions affect the variadic tail: float becomes double, and integer promotions apply to narrow integer types~\cite{wg14-2011-c11}.
The wrapper builds storage for the promoted values and uses the variadic preparation API with fixed and total argument counts.
System V AMD64 also uses the low byte of \code{RAX} to convey an upper bound on vector-register use~\cite{matz-2012-amd64}.
libffi supplies that mechanism when the call is described correctly.

Our test calls a variadic sum with fixed count $2$ and doubles $1.25$ and $2.5$.
The result is $3.75$, matching the compiler-generated direct call.
An unpromoted float descriptor in the variadic tail is rejected by the tested libffi implementation.
Declaring a variadic function as fixed-signature, even for one observed argument list, is not a general substitute.

\section{Exercises}
\label{sec:ch70-exercises}
\begin{enumerate}
\item\label{ex:ch70-eight} Assign eight uint64 arguments and evaluate $\sum_{i=1}^{8}i\cdot i$, the native test's weighted sum.
\item\label{ex:ch70-rollback} Explain why the last double in Table~\ref{tab:ch70-pressure} does not use XMM1.
\item\label{ex:ch70-hidden} A function returns a memory-class aggregate and takes six uint64 arguments.
How many explicit arguments remain in integer registers?
\item\label{ex:ch70-promote} A variadic tail supplies a float and signed char with values $1.25$ and $7$.
State the promoted types and why the original small objects are invalid slots for those types.
\end{enumerate}

\section{Further reading}
\label{sec:ch70-reading}
The AMD64 supplement defines the classification, register, stack, and variadic rules used here~\cite{matz-2012-amd64}.
The libffi manual specifies descriptor lifetimes, slot indirection, and widened returns~\cite{libffi-2026-manual}.
C11 distinguishes language promotions from ABI placement~\cite{wg14-2011-c11}.
